\section{Mastumoto-Amano Normal Form}

Single-qubit Clifford + T circuits can be written in the form of 
$$
	C_n T C_{n-1} T \cdots C_1 T C_0
$$
where each $C_n$ is some Clifford circuit.

Matsumoto and Amano \cite{matsumoto} showed that every single qubit Clifford + T circuit can be written in the following more restrictive form. 
$$
	W_n T W_{n-1} T \cdots W_1 T C
$$
where $W_n \in \{I, H, SH\}$, $W_i \in \{H, SH\}$ and $C$ is some Clifford circuit.
This is called the \textit{Matsumoto-Amano normal form}.

The importance of the Mastumoto-Amano normal form is summarised in the following theorem.

\begin{theorem}\label{thm:matsumoto-amano}
	T-count of the normal form is optimal.
\end{theorem}

We can, in addition, show that 
\begin{lemma}
\begin{enumerate}
		\item Every single qubit Clifford + T circuit can be written uniquely in the Matsumoto-Amano normal form. 
		\item A linear time algorithm exists to compute the normal form of a given single qubit Clifford + T circuit, with respect to the number of gates of the input.
	\end{enumerate}
\end{lemma}

\subsection{Existence}

We shall prove Theorem \ref{thm:matsumoto-amano} by the following lemmas.
Recall we have the following elementary quantum gates:
$$
H = \frac{1}{\sqrt{2}}\begin{pmatrix}
1 & 1 \\
1 & -1
\end{pmatrix}, \quad S = \begin{pmatrix}
1 & 0 \\
0 & i
\end{pmatrix}, \quad T = \begin{pmatrix}
1 & 0 \\
0 & e^{i\pi/4}
\end{pmatrix}, \quad \omega = e^{i\pi/4} I
$$

In all of the following lemmas, let $\mathcal{C}$ denote the one-qubit Clifford group generated by $H, S,$ and $\omega$. 
Let $\mathcal{L}$ denote the its subgroup generated by $S, X, $ and $\omega$, where $X = HS^2H$ is the Pauli X gate. 
Let $\C ' =  \C \backslash \L$.
Let $\mathcal{H} = \{I, S, SH\}$ and $\mathcal{H}' = \{H, SH\}$.

\begin{lemma}
	$\L$ has three left cosets in $\C$, namely $\L, H\L,$ and $SH\L$. 
	So $\C = \H \L$ and $\C' = \H' \L$.
\end{lemma}

\begin{proof}
	It is well known that the order of $\C$ is $192 = 64 \cdot 3$ \cite{mastel2025cliffordtheorynqubitclifford}. 
	The group $\L = (S, X, \omega)$ has the following structure
	$$
	XS = S^3 \omega^2 X, \omega S = S \omega, \omega X = X \omega, S^4 = I, X^2 = I, \omega^8 = I
	$$
	so all element of $L$ can be written as $\omega^i S^j X^k$ for some power $i, j, k$.
	As $\omega$ has order $8$, $S$ has order $4$ and $X$ has order $2$, $|L| = 4 \cdot 2 \cdot 8 = 64$.

	This means $L$ has three left cosets.
	Now, a simple calculation shows that $H\L$ and $SH\L$ are the other two left cosets, so $\C = \H \L$.
	$\C'$ is the complement of $\L$ in $\C$, so $\C' = \H' \L$.
\end{proof} 

\begin{lemma}
	$\L \H ' \subseteq \H ' \L$
\end{lemma}

\begin{proof}
	Since $S \in \L$, we have $\L\H' = \L S \cup \L SH \subset \L H$, which is non-trivial right coset of $\L'$.
	The non-trivial right coset must be a subset of the $\C' = \C \backslash \L = \H'\L$.
\end{proof}

\begin{lemma}
	$\L T = T \L$.
\end{lemma}

\begin{proof}
	$T$ commutes with generators of $\L$ in the following way
	$$
	ST = TS, XT = T XS \omega^{-1}, \omega T = T \omega 
	$$
	Thus for all $l \in \L$ we have another $l' \in \L$ s.t. $lT = Tl$, which means $\L T = T \L$.	
\end{proof}

\begin{lemma}\label{TLT=L}
	$T\L T = \L$
\end{lemma}

\begin{proof}
	Notice that $T\L T = TT \L = S\L = \L$, as $S \in \L$.	
\end{proof}

We are now ready to prove the existence of the Matsumoto-Amano normal form.

\begin{proof}[Proof: Uniqueness of Matsumoto-Amano normal form]
	By definition, Clifford + T circuits can be written as 
	$$M = C_n T C_{n-1} T \cdots C_1 T C_0.$$
	If any of the intermediate $C_i \in \L$ for $i\in \{1, 2, \cdots, n-1\}$, use $\ref{TLT=L}$ to reduce $TC_iT$ into another Clifford, which will merge the its adjacent $C_{i+1}, C_{i-1}$. 
	Doing so reduces $T$ counts by $2$.
	
	After reducing all intermediate $C_i$, our gates will be in the set 
	\begin{dmath}
	M 
	= \C T \C' T \cdots T \C' T \C
	=  \H \L T \H' \L T \cdots T \H' \L T \H \L
	\end{dmath}
	We can move all the $\L$ to the right of $\H'$ and $T$, this gives 
	$$
	M = \H T \H' T \cdots T \H' T \H \L
	$$
	which is the desired Matsumoto-Amano normal form.
\end{proof}

This proof is constructive, giving us a linear-time algorithm to reduce any Clifford + T circuit into the Matsumoto-Amano normal form, which is the second part of Theorem \ref{thm:matsumoto-amano}.

\begin{proof}
	Since $\C$ is a finite group, we may assume computation in $\C$ takes constant time using a hash table.
	As $\L$ is a subgroup of $\C$, identifying whether an element is in $\L$ also takes constant time.
	
	Starting with a given input in the form
	$$M = C_n T C_{n-1} T \cdots C_1 T C_0.$$
	we can check from left to right whether $C_i \in \L$ for $i\in \{1, 2, \cdots, n-1\}$. 
	If $C_i$ is in $\L$, replace $C_{i+1} T C_i T C_{i-1}$ with an element in $\L$. 
	At each step, we are moving rightwards by $2$ gates, so this process takes at most $O(n)$ time.

	After reducing all intermediate $C_i$, we can move all $\L$ to the right of $\H'$ and $T$ with the same process.
\end{proof}

\subsection{Optimality of T-count}

The aim of this section is to show the following theorem.

\begin{theorem}\label{thm:matsumoto-amano-uniqueness}
	Two Matsumoto-Amano normal forms represent the same unitary if and only if they are identical as circuits. 
	That is, they can be written down in the exact same way.
\end{theorem}

Assuming the correctness of theorem \ref{thm:matsumoto-amano-uniqueness}, we can show that the $T$-count of the Matsumoto-Amano normal form is minimal.

\begin{corollary}
	The number of $T$ gates in the Matsumoto-Amano normal form is minimal among all Clifford + T circuits representing the same unitary.
\end{corollary}

\begin{proof}
	Let $M$ be a Clifford + $T$ unitary with the minimal number of $T$ gates in its circuit representation. 
	Such a minimal representation exists by Zorn's lemma.

	We can reduce $M$ into its Matsumoto-Amano normal form $M'$ by the procedure described in the proof of the existence. 
	In that procedure, the number of $T$ gates may only decrease, so $M'$ will have the same number of $T$ gates as $M$.

	By uniqueness, all representations of $M$ can be reduced to the same normal form.
	This concludes the proof.
\end{proof}

The proof of \autoref{thm:matsumoto-amano-uniqueness} requires more sophisticated analysis on the structure of $\C$ and $\L$, which we will achieve in the following paragraphs using more tools in algebra.

As we have shown in  \autoref{sec:bloch-sphere}, all one-qubit unitaries, up to an insignificant global phase, can be represented as a rotation on the Bloch sphere.
As rotations in 3 dimensions are represented by $SO(3)$, we can represent each one-qubit unitary as an element in $SO(3)$, which is called the \textit{Bloch sphere representation} of the unitary.

\begin{definition}[Bloch Sphere Representation]
	Let $U SU(2)$, denote $\hat{U} \in SO(3)$ as the rotation on the Bloch sphere induced by $H$.
\end{definition}

For example, the Bloch sphere representation of $H$, $S$, and $T$ are:
\begin{equation}\label{eq:bloch-representation}
\hat{H} = \begin{pmatrix}
0 & 0 & 1 \\
0 & -1 & 0 \\
1 & 0 & 0
\end{pmatrix},  
\hat{S} = \begin{pmatrix}
0 & -1 & 0 \\
1 & 0 & 0 \\
0 & 0 & 1
\end{pmatrix},  
\hat{T} = \frac{1}{\sqrt{2}}\begin{pmatrix}
1 & -1 & 0 \\
1 & 1 & 0 \\
0 & 0 & \sqrt{2} \\
\end{pmatrix}
\end{equation}

\begin{remark}\label{remark:clifford-so3}
	Let $\hat{\C}$ denote the Bloch sphere representation of Clifford group, which is generated by $\hat{H}$ and $\hat{S}$.
	$\hat{H}$ is the linear transformation that permutes the $x, y, z$ axis to $z, -y, x$, and $\hat{S}$ permutes the $x, y, z$ axis to $-y, x, z$.
	Thus, all elements of  $\hat{\C}$ must be some permutation of the $x, y, z$ axis to $\pm x, \pm y, \pm z$.

	Let us use a temporary notation $(a, b, c)$ to denote the linear map that sends $x, y, z$ to $a, b, c$ respectively.

	Note 
	\begin{align*}
		\hat{H} &= (z, -y, x) \\
		\hat{S} \hat{H} &= (z, x, y)\\ 
		\hat{S}  &= (y, -x, z) \\
		\hat{S}^3 &= (-y, x, z)
	\end{align*}

	Both of $\hat{H}$ and $\hat{S}$ are odd permutations, which also brings one negative sign. 
	Ignoring the signs, we $\hat{S} \hat{H}$ and $\hat{S}$ are exactly the generators of the symmetric group $S_3$. 
	Combining these two observations, we see $\hat{C}$ consists of odd permutations of $x, y, z$ with an odd number of negative signs and even permutations of $x, y, z$ with an even number of negative signs.

	It turns out this is the only rule that needs to be satisfied for a linear map to be in $\hat{\C}$.
	This is because $\hat{S}$ swaps the first two axes while negating the second, and $\hat{S}^3$ swaps the first two axes while negating the first.
	Paired with $\hat{S}\hat{T}$, they can generate any permutation of two axes while negating either one of them.

	This means $\hat{\C}$ consists of $3! \cdot 8 / 2 = 24$ elements, which are exactly the elements of $SO(3)$ with entries only in $\pm 1$.
\end{remark}

Remark \ref{remark:clifford-so3} states that all Cliffords have Bloch sphere representation with entries $\pm 1$. 
Equation \ref{eq:bloch-representation} shows that $T$ has Bloch sphere representation with entries that are the sum and product of integers and $\sq$. This is exactly the ring of numbers $\Z[\sq]$.

Since Clifford + T circuits are generated by Clifford and $T$, we have the following proposition.

\begin{proposition}
	The Bloch sphere representation of any Clifford + T circuit is an a matrix in $SO(3, \mathbb{Z}[\frac{1}{\sqrt{2}}])$.
\end{proposition}

We define a map $q: SO(3, \Z[\sq]) \rightarrow Mat_{3, 3}(\Z / 2\Z)$ in the following way. 
For any $u \in SO(3, \Z[\sq])$, 
\begin{enumerate}
	\item Find the smallest $k$ such that $\left(\sq\right)^k u$ has entries in $\Z[\sqrt{2}]$.
	\item each entry of $\left(\sq\right)^k u$ can be written as $a + b\sqrt{2}$ for some integers $a, b$, we define the corresponding entry of $q(u)$ to be $a \mod 2$.
\end{enumerate}

The number $k$ is the \textit{least denominator exponent} of $u$.
We call $p(u)$ the \textit{parity form} of $u$.

Define an equivalence relationship $\sim$ on $SO(3, \Z[\sq])$ by $u \sim v$ if and only if $u $ and $v$ differ by a permutation of columns.

We have the following important lemma.

\begin{lemma}
	Let $M$ be the Bloch sphere representation of a Matsumoto-Amano normal form with least denominator exponent $k$.
	Then exactly one of the following holds:
	\begin{enumerate}
		\item $k = 0$, and $M$ is Clifford. 
		\item $k > 0$, and $q(M) \sim \begin{pmatrix}
1 & 1 & 0 \\ 
1 & 1 & 0 \\ 
0 & 0 & 0
\end{pmatrix}$ if the first gate of the normal form is $T$
		\item $k > 0$, and $q(M) \sim \begin{pmatrix}
0 & 0 & 0 \\ 
1 & 1 & 0 \\ 
1 & 1 & 0
\end{pmatrix}$ if the first gate of the normal form is $HT$
		\item $k > 0$, and $q(M) \sim \begin{pmatrix}
1 & 1 & 0 \\ 
0 & 0 & 0 \\ 
1 & 1 & 0
\end{pmatrix}$ if the first gate of the normal form is $SHT$
	\end{enumerate}
\end{lemma}

\begin{proof}
	Any Matsumoto-Amano normal form can be written as some Clifford left multiplied by some sequence of $HT$ and $SHT$, which terminates with an optional $T$.

	Elements of the Clifford group have a parity form of identity. 
	Direct computation shows that, when left multiplied by $T$, $HT$, or $SHT$, the parity form changes as illustrated by the scheme below, where $k++$ denotes the increment of the least denominator exponent by $1$.
	From which it is clear that the parity form of a normal form is completely determined by the first gate of the normal form.
\ctikzfig{ma-form}

\end{proof}

We are now ready to prove the uniqueness of the Matsumoto-Amano normal form.

\begin{proof}[Proof: Uniqueness of Matsumoto-Amano normal form]
	Let $M$ and $N$ be two Matsumoto-Amano normal of the same unitary map.
	Without loss of generality, let $T$ gate counts of $M$ be smaller than or equal to those of $N$.

	Since $M$ and $N$ represent the same unitary, they must have the same parity form of Bloch sphere representation. This means the starting gate of $M$ and $N$ must be the same, which is either $T$, $HT$, or $SHT$.
	We can then cancel the first gate of $M$ and $N$, and repeat the
	process until we have cancelled all $T$ gates of $M$, leaving us with a Clifford circuit with parity form of identity.
	At this moment, the parity form of $N$ must also be identity, so $N$ is also left with a Clifford circuit equal to that of $M$.
\end{proof}

